\documentclass{include/thesisclass}
% Main File - Based on thesisclass.cls
% Comments are mostly in English
% ------------------------------------------------------------------------------
% Further files in folder:
%  - include/cmds.tex (for macros and additional commands)
%  - include/kitlogo.pdf (for titlepage)
%  - lit.bib (bibtex bibliography database)
%  - include/titlepage.tex (for layout of titelpage)
% ------------------------------------------------------------------------------
% Useful Supplied Packages:
% amsmath, amssymb, mathtools, bbm, upgreek, nicefrac,
% siunitx, varioref, booktabs, graphicx, tikz, multicol





%% -------------------------
%% |    Thesis Settings    |
%% -------------------------
% english or ngerman (new german für neue deutsche Rechtschreibung statt german)
\SelectLanguage{english}
% details on this thesis
\newcommand{\thesisauthor}{Lukas Powalla}
\newcommand{\thesistopic}{Cloaced grids on solar cells}
\newcommand{\thesisentopic}{Getarnte Kontakte auf Solarzellen}
\newcommand{\thesislongtopic}{Very long and very detailed description of the very interesting thesis topic (only necessary for pdfsubject tag).}
\newcommand{\thesisinstitute}{Institut für Angewandte Physik}
\newcommand{\thesisreviewerone}{Prof. Dr. Martin Wegener}
\newcommand{\thesisreviewertwo}{Martin Schumann}
\newcommand{\thesisadvisorone}{} % to use: enter names and uncomment in titlepg
\newcommand{\thesisadvisortwo}{}
\newcommand{\thesistimestart}{01.02.2016} % on titlepage
\newcommand{\thesistimeend}{01.05.2016} % on titlepage
\newcommand{\thesistimehandin}{20.04.2016} % on second page 'preamble'
\newcommand{\thesispagehead}{Bachelor Thesis: \thesisentopic} % page heading





%% ---------------------
%% |    PDF - Setup    |
%% ---------------------
% This information will appear embed into the PDF file as meta data, but will 
% not be printed anywhere
\hypersetup
{
    pdfauthor={\thesisauthor},
    pdftitle={Bachelorarbeit: \thesistopic},
    pdfsubject={\thesislongtopic},
    pdfkeywords={kit,physik,bachelor,thesis,\thesisauthor}
}





%% --------------------------------------
%% |    Settings for Word Separation    |
%% --------------------------------------
% Help for separation:
% In German package the following hints are additionally available:
% "- = Additional separation
% "| = Suppress ligation and possible separation (e.g. Schaf"|fell)
% "~ = Hyphenation without separation (e.g. bergauf und "~ab)
% "= = Hyphenation with separation before and after
% "" = Separation without a hyphenation (e.g. und/""oder)

% Describe separation hints here:
\hyphenation
{
    über-nom-me-nen an-ge-ge-be-nen
    %Pro-to-koll-in-stan-zen
    %Ma-na-ge-ment  Netz-werk-ele-men-ten
    %Netz-werk Netz-werk-re-ser-vie-rung
    %Netz-werk-adap-ter Fein-ju-stier-ung
    %Da-ten-strom-spe-zi-fi-ka-tion Pa-ket-rumpf
    %Kon-troll-in-stanz
}





%% -----------------------
%% |    Main Document    |
%% -----------------------
\begin{document}
    % Titlepage and ToC
    \FrontMatter

    % coordinates for background border
\newcommand{\diameter}{20}
\newcommand{\xone}{-15}
\newcommand{\xtwo}{160}
\newcommand{\yone}{15}
\newcommand{\ytwo}{-253}




\begin{titlepage}
    % background border
    \begin{tikzpicture}[overlay]
    \draw[color=gray]
            (\xone mm, \yone mm)
      -- (\xtwo mm, \yone mm)
    arc (90:0:\diameter pt)
      -- (\xtwo mm + \diameter pt , \ytwo mm)
        -- (\xone mm + \diameter pt , \ytwo mm)
    arc (270:180:\diameter pt)
        -- (\xone mm, \yone mm);
    \end{tikzpicture}



    % KIT image and sign for faculty of physics
    \begin{textblock}{10}[0,0](4.5,2.5)
        \includegraphics[width=.25\textwidth]{include/kitlogo.pdf}
    \end{textblock}
    \changefont{phv}{m}{n}    % helvetica
    \begin{textblock}{10}[0,0](5.5,2.2)
        \begin{flushright}
            \Large FAKULTÄT FÜR PHYSIK\\\thesisinstitute
        \end{flushright}
    \end{textblock}



    % horizontal line
    \begin{textblock}{10}[0,0](4.2,3.1)
        \begin{tikzpicture}[overlay]
        \draw[color=gray]
                (\xone mm + 5 mm, -12 mm)
          -- (\xtwo mm + \diameter pt - 5 mm, -12 mm);
        \end{tikzpicture}
    \end{textblock}



    % begin of text part
    \changefont{phv}{m}{n}    % helvetica
    \centering



    % thesis topic (en and ge)
    \vspace*{3cm}
    \Huge\thesistopic\\
    \huge(\thesisentopic)\\



    % author name and institute
    \vspace*{2cm}
    \Large Bachelorarbeit\\von\\
    \vspace*{1cm}
    \huge\thesisauthor\\
    \vspace*{1cm}
    \Large am \thesisinstitute



    % possible frontimage - thanks to JabberWok
    % for publishing the img under GNU Document License
    \vspace*{1.5cm}
    \includegraphics[scale=0.7]{./include/frontimage1.jpg}\\



    % examiners (Referenten)
    \vspace*{1.5cm}
    \Large
    \begin{center}
        \begin{tabular}[ht]{l c l}
        \iflanguage{english}{professor}{Referent}: 
            & \hfill & \thesisreviewerone\\
        \iflanguage{english}{adviser}{Korreferent}: 
            & \hfill & \thesisreviewertwo\\
        % uncomment if you want to provide info on your advisors
        %\iflanguage{english}{Advisor}{Betreuender Mitarbeiter}: 
        %    & \hfill & \thesisadvisorone\\
        %\iflanguage{english}{Second Advisor}{Zweiter betreuender Mitarbeiter}: 
        %    & \hfill & \thesisadvisortwo\\
        \end{tabular}
    \end{center}



    % working time
    \vspace{1cm}
    \begin{center}
        \large{{Bearbeitungszeit}: \thesistimestart \hspace*{0.25cm} -- %
                                   \hspace*{0.25cm} \thesistimeend}
    \end{center}



    % lowest text blocks concerning the KIT
    \begin{textblock}{10}[0,0](4,16.8)
        \tiny{KIT -- Universität des Landes Baden-Württemberg und nationales 
             }
    \end{textblock}
    \begin{textblock}{10}[0,0](14,16.75)
        \large{\textbf{www.kit.edu}}
    \end{textblock}
\end{titlepage}

    \chapter*{Erklärung zur Selbstständigkeit}
Ich versichere, dass ich diese Arbeit selbstständig verfasst habe und keine %
anderen als die angegebenen Quellen und Hilfsmittel benutzt habe, die %
wörtlich oder inhaltlich übernommenen Stellen als solche kenntlich gemacht und %
die Satzung des KIT zur Sicherung guter wissenschaftlicher Praxis in der %
gültigen Fassung vom 17.05.2010 beachtet habe.\\

\vspace{1cm}

\renewcommand{\arraystretch}{0} % for spacing in the tabular environment

\begin{flushright}
	\begin{tabular}{rr}
		Karlsruhe, den \thesistimehandin, & \hspace*{5cm}\\[0mm]
		\cline{2-2}\\[2mm]    % the last line has height 2mm due
		& \thesisauthor       % to \arraystretch=0
	\end{tabular}
\end{flushright}

\vfill

\begin{flushright}
	Als Ansichtsexemplar genehmigt von\\
	\vspace{1cm}
	\begin{tabular}{rr}
		Karlsruhe, den \thesistimehandin, & \hspace*{5cm}\\[0mm]
		\cline{2-2}\\[2mm]    % the last line has height 2mm due
		& \thesisreviewerone  % to \arraystretch=0
	\end{tabular}
\end{flushright}

\renewcommand{\arraystretch}{1}

\cleardoublepage


    \begingroup \let\clearpage\relax    % in order to avoid listoffigures and
    \tableofcontents                    % listoftables on new pages
    %\listoffigures
    %\listoftables
    \endgroup
    %\cleardoublepage



    % Contents
    \MainMatter

    \chapter{Introduction}
    During the last few years, the efficiency of solar cells could be raised significantly through research and development on different aspects of the solar cell. Besides reducing electrical losses and improving the materials used in solar cells (such as the material for the semi-conductor), we can also try to improve optical properties of the solar cell. On the one hand-side, it is important to reduce reflections and on the other hand-side, it is important to guide as much light as possible to the optical active areas. 
In this Bachelorthesis, we will concentrate on guiding as much light as possible to the optical active areas.
One issue concerning this is cloaking the contact fingers and bus-bars on solar cells to guide the light, which would have hit the contact grid on solar cells, to the optical active areas. This can be used to enhance the efficiency of solar cells. 


	\chapter{Description of the Problem}
    \section{state of the art}
In order to improve the optical properties of solar cells a lot of research has been done. It is important to minimise reflections on the sunny side of the solar cell and in addition to that it is also important to guide as much light as possible to the optical active areas of the solar cell. 
Reflected light and light that hits the contact grid and the busbar is lost.

The question I want to deal with in my Bachelorthesis is how to minimise light that hits the contact grid and the busbar. You can reach this by either decreasing the contact grid and the busbars. This can be reached through using back contact solar cells \cite{kerschaver2006back} or emitter wrap through solar cells \cite{gee1992emitter}. 

Another approach is to design the sun-side of the solar cell in a way that the light in the end hits the optical active area. 


    \chapter{Simulations for continuous surface design}    
    \section{blending of the two one dimensional solutions}

\section{Simulation rectangular unit cell}

\section{Simulation squared unit cell}
	\chapter{Simulations for Fresnel design}
	\section{Construction of the Fresnel design}

\section{Simulation for rectangular unit cell}

\section{Simulation for squared unit cell}
	\chapter{Simulations for meta-surface design }
	\section{construction of the meta-surafce design}

\section{Simulation for rectangular unit cell}

\section{Simulation for squared unit cell}


    \chapter{Conclusions}
    \input{./chap/conclusion.tex}

    % appendix for more or less interesting calculations
    \Appendix
    \chapter*{\appendixname} \addcontentsline{toc}{chapter}{\appendixname}
    % to make the appendix appear in ToC without number. \appendixname = 
    % Appendix or Anhang (depending on chosen language)
    \section{First Appendix Section}
Wonderful Appendix!

 %\cleardoublepage


    
    \TheBibliography
    % BIBTEX
    % use if you want citations to appear even if they are not referenced to: 
    % \nocite{*} or maybe \nocite{Kon64,And59} for specific entries
    \nocite{*}
    \bibliographystyle{babalpha}
    \bibliography{lit}
    % THEBIBLIOGRAPHY
%    \begin{thebibliography}{000}
%        \bibitem{ident}Entry into Bibliography.
%    \end{thebibliography}
\end{document}
